\clearpage
\pdfbookmark[1]{Zum Gebrauch des Komparators}{using-the-comparator}
\begin{center}\textbf{Zum Gebrauch des Komparators}\end{center}
\noindent Diese mit KI erstellte Fassung in modernem Deutsch beruht auf dem deutschen Original \textit{Zum Gebrauch des Comparators} in Carl Friedrich Gauss, \textit{Werke}, Band~V, S.~632--635, herausgegeben von der Königlichen Gesellschaft der Wissenschaften zu Göttingen. Siehe das \href{https://gdz.sub.uni-goettingen.de/id/PPN236006339}{Göttinger Digitalisat} und den \href{https://archive.org/download/Werkecarlf05gausrich/page/n643}{zweiten Abdruck (1877)}.
\par\smallskip\noindent Zur messtechnischen Terminologie siehe John S. Beers und William B. Penzes, \href{https://nvlpubs.nist.gov/nistpubs/jres/104/3/html/j43bee.htm#c1pz3}{\textit{The NIST Length Scale Interferometer}}, \textit{Journal of Research of the National Institute of Standards and Technology} 104(3), Mai--Juni 1999, §§1.3--1.4, und die Physikalisch-Technische Bundesanstalt, \href{https://oar.ptb.de/files/download/460.2016.pdf}{\textit{Jahresbericht 2016}}, S.~18--19.
\clearpage
